\lecture{6}
\title{TM Variants, Church-Turing Thesis}

\begin{document}

Grades for PSET 1 have been released on Gradescope.
We have a quiz next Tuesday (October 6th) from 7:30 PM to 9:00 PM in Walker;
sample questions will be posted on the course website.

\subsection{Turing Machines (Recap) \& The Church-Turing Thesis}

Recall the definition of a Turing machine from
\href{/18.404/lectures/5/#turing-machines}{Lecture 5}.
Importantly, for any word $w$, a Turing machine $M$ can either:
\begin{center}
    \texttt{ACCEPT} ~ (enter $q_{\mathrm{acc}}$)
    ~~~ $\parallel$ ~~~
    \texttt{REJECT BY HALTING} ~ (enter $q_{\mathrm{rej}}$)
    ~~~ $\parallel$ ~~~
    \texttt{REJECT BY LOOPING} ~ (never halt)
\end{center}
The Turing machine will be our model for a general-purpose computer.

\textbf{Question.}
Why is a one-tape TM our choice of model?
Why not anything else? \\
\textcolor{grey}{(Recall that a PDA with two stacks is strictly more powerful
than a PDA with just one stack.)}

\emph{Answer:}
As we'll eventually show in this course,
it turns out that a one-tape TM
is equally as powerful as, say, Python or Java,
because one can compile code in any one of these ``languages''
into any other.  In fact\dots

\begin{theorem}{The Church-Turing Thesis}
    Any function computable by an ``algorithm''
    is computable by a Turing machine.
\end{theorem}

(The Church-Turing thesis is more of a heuristic than a theorem.)

\textbf{Definition.}
Also recall from \href{/18.404/lectures/5/#turing-machines}{Lecture 5}
the following properties of languages and TMs.
\begin{itemize}
    \item A Turing machine $M$ is a \textbf{decider}
    if it halts on every input.
    \item A language $A$ is \textbf{Turing-recognizable}
    if $A = L(M)$ for some Turing machine $M$.
    \item A language $A$ is \textbf{decidable}
    if $A = L(M)$ for some \emph{decider} $M$.
\end{itemize}

\subsection{Multi-Tape Turing Machines}

\textbf{Definition.}
We say that a \textbf{multitape Turing machine} $M$
is a TM with multiple tapes, naturally.
\begin{center}
    \frame{\includegraphics[width=0.5\textwidth]{lecture-06/multitape-TM}}
\end{center}
In particular, the transition function of
a multitape Turing machine with $k$ tapes
has type $\delta \colon Q \times \Gamma^k \to
Q \times \Gamma^k \times \{\mathrm{L}, \mathrm{R}\}^k$.

\begin{theorem}{Multitape = Single Tape}
    A language $A$ is Turing-recognizable
    if and only if $A = L(M)$ for some \emph{multitape} TM $M$.
\end{theorem}

\begin{proof}
    The $\implies$ direction is obvious.
    For the $\impliedby$ direction, a one-tape TM $S$
    can simulate a multitape TM $M$ like so.
    \begin{center}
        \frame{\includegraphics[width=0.9\textwidth]{lecture-06/multitape-TM-proof}}
    \end{center}
    The idea is to simply store the concatenation of all $k$ tapes
    in a single tape.
    \begin{itemize}
        \item Separate each of the $k$ tapes with a \str{\#} symbol.
        \begin{itemize}
            \item If one of the $k$ tapes writes in a blank space---thereby
            ``extending its length''---then we simply
            shift everything to the right.
        \end{itemize}
        \item Keep track of the head position on each of the $k$ tapes
        by writing $\str{c}^{\bullet}$ if the head is pointing to $\str{c}$.
    \end{itemize}
    Details omitted, but they're not hard to implement. ~ $\blacksquare$
\end{proof}

\begin{remark}
    It follows that ``a PDA with 100 stacks plus a tape'' is equally as powerful
    as ``an FA plus a tape''.
\end{remark}

\subsection{Nondeterministic Turing Machines}

\textbf{Definition.}
We say that a \textbf{nondeterministic Turing machine (NTM)} $M$ is a TM
that's nondeterministic, naturally.
The transition function has type
$\delta \colon Q \times \Gamma \to
\mathcal{P}(Q \times \Gamma \times \{\mathrm{L}, \mathrm{R}\})$.

\begin{theorem}{NTM = TM}
    A language $A$ is Turing-recognizable if and only if
    $A = L(N)$ for some NTM $N$.
\end{theorem}

\begin{proof}
    The $\implies$ direction is obvious.

    For the $\impliedby$ direction, the idea is for a one-tape TM $M$
    to BFS through the tree of computation paths of $N$.
    \begin{center}
        \frame{\includegraphics[width=0.5\textwidth]{lecture-06/BFS}}
    \end{center}
    In the single tape of $M$,
    we simply store the state of all \emph{threads} of $N$,
    moving down one layer at a time.
    \begin{center}
        \frame{\includegraphics[width=0.75\textwidth]{lecture-06/NTM-proof}}
    \end{center}
    As before, threads are separated with $\str{\#}$,
    and tape heads are indicated by the $\bullet$ in $\str{c}^{\bullet}$.
    Furthermore, each block begins with a cell that stores
    the FA state that the thread lands at (e.g., $q_2$, $q_3$, or $q_5$,
    as shown above).

    So the TM should BFS through all threads layer-by-layer
    and return \str{ACCEPT} if any thread returns \str{ACCEPT}.
    (Details omitted.)
    ~ $\blacksquare$
\end{proof}

\subsection{Enumerators (Turing Machines with Printers)}

\textbf{Definition.}
We say that an \textbf{enumerator} $E$ is a Turing machine
that takes no inputs and prints outputs.  Specifically,
\begin{itemize}
    \item Its workspace tape is initially completely empty.
    \item At any point, $E$ can move to a \str{print} state
    and print the contents of the workspace.
\end{itemize}
We let $L(E)$ denote the set of strings that $E$ eventually prints.

\begin{theorem}{Enumerator = TM}
    A language $A$ is Turing-recognizable if and only if
    $A = L(E)$ for some enumerator $E$.
\end{theorem}

\begin{proof}
    We first consider the $\impliedby$ direction:
    given any $A$ for which $A = L(E)$,
    we construct a TM $M$ for which $A = L(M)$.
    \begin{quote}
        \emph{Solution:}
        On input $w$, simply simulate $E$ on a blank tape.
        During the simulation, check everything it prints for
        whether it matches $w$.
        If there's ever a match, return \str{ACCEPT}.
        ~ $\square$
    \end{quote}
    Now for the $\implies$ direction:
    given any $M$ for which $A = L(M)$,
    we construct an enumerator $E$ that prints out $A$.
    \begin{quote}
        \emph{Solution:}
        The rough idea is brute force: for each $x \in \Sigma^*$,
        simulate $M$ on $x$, and print $x$ if $M$ accepts.
        There are \emph{two} issues:
        \begin{itemize}
            \item If we try to simulate one $x \in \Sigma^*$ at a time,
            then we might fail if a simulation of $M$ on $x$ loops.
            \item If we try to simulate one layer at a time,
            then we'll never get past the first layer because
            $\Sigma^*$ is infinite.
        \end{itemize}
        (In a rough sense, the former is DFS, and the latter is BFS---but
        this time, they both fail.)

        The fix is similar to how we prove $|\mathbb{Q}| = |\mathbb{N}|$;
        do a ``zig-zag'' across both $\Sigma^*$ and layers, like so.
        \begin{center}
            \frame{\includegraphics[width=0.4\textwidth]{lecture-06/enumerator-proof}}
        \end{center}
        More formally, in the $i^{\text{th}}$ step,
        run $(x_1, x_2, \dots, x_i)$ up to the $i^{\text{th}}$ layer,
        and do this for all $i \in \mathbb{N}$. ~ $\square$
    \end{quote}
    This completes the proof of both directions.
    (Again, many details omitted.) ~ $\blacksquare$
\end{proof}

\end{document}